\section{Comparison with OCO-2}\label{sec:oco2}

The companion paper \cite{nmkc} applies the same construction to the OCO-2 radiative-transfer
emulation problem of Lamminp\"a\"a et al.~\cite{lamminpaa}: per spectral band (O$_2$, weak
CO$_2$, strong CO$_2$), a reduced atmospheric state of twenty or twenty-four coordinates maps to
forty reduced radiance coefficients, with errors on the standardized coefficients and on the
reconstructed radiance, under the ten-split protocol of Section~\ref{sec:data}. It gives a
comparison at a higher input dimension.

There a kernel on the state trails the network by a factor of nine, four and five on the reduced
coefficients: the isotropic Mat\'ern reaches $40.6\%$, $58.6\%$ and $43.4\%$ on the three bands
against $4.4\%$, $16.4\%$ and $8.2\%$ for the network. Four ways of learning the kernel's metric
on the state were run on the same ten splits on all three bands: a kernel-flow metric
\cite{owhadiyoo}, a Mahalanobis metric under the same flow, the gradient outer product of the
recursive feature machine \cite{rfm}, and per-input scales fitted by empirical Bayes. The last is
the best of the four on every band, taking the isotropic kernel's $40.6\%$, $58.6\%$ and $44.3\%$
on the coefficients to $18.7\%$, $25.3\%$ and $16.9\%$, and on O$_2$ its radiance from $0.233\%$ to
$0.150\%$; the two flow metrics agree within half a point of each other at $26.9\%$, $29.8\%$ and
$23.9\%$, the recursive feature machine reads $28.2\%$, $47.9\%$ and $31.3\%$, and an additive
kernel $38.2\%$, $46.9\%$ and $47.7\%$. (The isotropic kernel on strong CO$_2$ reads $44.3\%$ here
and $43.4\%$ in the baseline experiment. The metric experiment also chooses the Mat\'ern smoothness
on validation and takes $\nu=3/2$ on this band at every split; the baseline fixes $\nu=5/2$, and its
choice of scale and nugget changes from split to split, six splits landing near $41\%$, three near
$48\%$ and one at $44\%$, hence its spread of $\pm3.4$.)
The feature kernel remains more accurate, by a margin that depends on the band: the exact Mat\'ern
regression on the network's frozen features reads $4.1\%$, $16.3\%$ and $8.1\%$ on the same splits,
so the best state metric is $4.6$ times that on O$_2$ but only $1.6$ and $2.1$ times it on the two
CO$_2$ bands. The empirical-Bayes metric is bimodal across splits on O$_2$, eight of ten landing near
$17.2\%$ and two at $24.7\%$, so its spread of $\pm3.2$ is thirty times the isotropic kernel's
$\pm0.1$. The exact Mat\'ern regression on the network's frozen features beats a ridge readout
refitted on the same features on every split, every band and every training loss, and a
per-coefficient choice among heads trained under different losses has lower mean error than the
release's own emulator under both reported metrics: it wins the reduced-coefficient comparison on
all ten splits on each band, and the radiance comparison on ten, nine and ten splits on O$_2$, weak
CO$_2$ and strong CO$_2$.

\section{Ten further emulation configurations}\label{sec:corpora}

Table~\ref{tab:corpora} extends the comparison to ten further configurations, each with
validation-based selection; the comparison concerns the families under a common protocol, not the
best method available for each problem.

The corpora are the advection equation with discontinuous initial data from the operator suite
of \cite{dehoop}, which is also the source of the structural-mechanics problem of \cite{nmkc};
Burgers' equation at three viscosities from PDEBench \cite{pdebench}; three one-step maps from
The Well \cite{thewell}: a two-dimensional turbulent radiative layer, an active-matter system and
a Helmholtz staircase; the ClimSim atmospheric physics emulation task \cite{climsim}; and the
paired correction-coefficient corpus of the physics-guided radiative-transfer work of \cite{pkan},
a third radiative-transfer emulation problem and the largest here, with about 488{,}000 training
rows. That corpus
is run in both of its configurations: from the atmospheric state alone, and with the low-fidelity
6S coefficients as additional inputs, which is a different and easier problem.

Advection is the 200-point input/output pair of the operator suite at full resolution. Burgers is
PDEBench's \texttt{1D\_Burgers\_Sols\_Nu}$\nu$ release, the map from the initial condition to the
solution at $t = 1$ on a grid subsampled by four to 256 cells, with the input used raw and the output
reduced to 128 principal coefficients. The three Well datasets are one-step maps of the full
multi-channel state with 256 principal coefficients in and out: the turbulent radiative layer on
its $128 \times 384$ grid (density, pressure, velocity), active matter on $256 \times 256$
(concentration, velocity and the two tensor fields), and the Helmholtz staircase on its two pressure
channels. ClimSim is the low-resolution release, 124 inputs to 128 tendencies, with 100{,}000
training rows drawn from the 10.1 million available and validation and test taken from the release's
own validation and scoring files. The correction-coefficient corpus is the paired 6S/libRadtran
Sentinel-2 set \cite{vermote,mayer}: wavelength, solar zenith, view zenith, relative azimuth,
aerosol optical depth, water vapor, ozone and elevation map to three libRadtran coefficients, split
by atmospheric state so that all bands of a state stay together; the multi-fidelity configuration
adds the three 6S coefficients, for input dimensions eight and eleven. For these two configurations
the error is
\begin{equation}\label{eq:coefficient-floor}
 E_{0.05}=\frac{1}{N}\sum_{i=1}^{N}
 \frac{\|Y_i-\widehat Y_i\|_2}{\max\{\|Y_i\|_2,0.05\}},
\end{equation}
and the same floor enters validation selection and the kernel-flow objective, which limits the
weight of coefficient vectors of small norm.

The families are those of Section~\ref{sec:methods} with two changes forced by scale. The
reference is a ridge regression linear in the inputs, since a cubic expansion of a 256-dimensional
input is not tractable. Kernel solves are exact when the training block has at most 20{,}000 rows
and use a Nystr\"om approximation \cite{nystrom} on 6{,}000 uniformly drawn landmarks above that,
which applies to ClimSim, the Helmholtz staircase at 20{,}384 rows and the two
correction-coefficient configurations. The convex stack and the per-coordinate selection are formed
from the kernels, the networks, the residual correction and the feature kernel; the ridge is a
reference row, not one of their members. Advection, Burgers, ClimSim and the two
correction-coefficient configurations have three to five random splits. The Well's releases carry
their own training, validation and test files, which are used unchanged, so the three
turbulent-radiative-layer runs share one split and differ only in network initialization (their
spread, of order $0.01$ points, measures initialization), and the active-matter and
Helmholtz-staircase rows are single runs.

\begin{table}[t]
\centering
% Generated by code/format_publication_tables.py; numerical cells preserved.
\textbf{A. Configurations}\par\medskip
{\fontsize{9}{11}\selectfont
\setlength{\tabcolsep}{4pt}
\begin{tabular}{lrrrr}
\toprule
corpus & $n$ & $d$ & $q$ & runs \\
\midrule
Advection & 19,000 & 200 & 200 & 5 \\
Burgers $\nu=0.1$ & 7,200 & 256 & 128 & 3 \\
Burgers $\nu=0.01$ & 7,200 & 256 & 128 & 3 \\
Burgers $\nu=0.001$ & 7,200 & 256 & 128 & 3 \\
Radiative layer & 7,200 & 256 & 256 & 3 \\
Active matter & 14,000 & 256 & 256 & 1 \\
Helmholtz staircase & 20,384 & 256 & 256 & 1 \\
ClimSim & 100,000 & 124 & 128 & 5 \\
RT coefficients, state & 487,773--487,823 & 8 & 3 & 5 \\
RT coefficients, +6S & 487,773--487,823 & 11 & 3 & 5 \\
\bottomrule
\end{tabular}}
\par\medskip\textbf{B. Reference, input kernels and network}\par\medskip
{\fontsize{9}{11}\selectfont
\setlength{\tabcolsep}{3pt}
\begin{tabular}{lrrrr}
\toprule
corpus & linear ridge & KRR & \thead{learned-metric\\KRR} & network \\
\midrule
Advection & 11.29$\pm$0.04 & 11.36$\pm$0.04 & 11.36$\pm$0.04 & 11.69$\pm$0.07 \\
Burgers $\nu=0.1$ & 13.89$\pm$0.53 & 2.81$\pm$0.28 & 2.90$\pm$0.33 & 20.34$\pm$0.61 \\
Burgers $\nu=0.01$ & 33.68$\pm$0.28 & 23.40$\pm$0.54 & 23.58$\pm$0.50 & 36.65$\pm$1.95 \\
Burgers $\nu=0.001$ & 38.38$\pm$0.12 & 31.66$\pm$0.16 & 33.23$\pm$1.73 & 39.98$\pm$0.65 \\
Radiative layer & 31.73 & 30.31 & 30.26$\pm$0.01 & 32.17$\pm$0.02 \\
Active matter & 37.24 & 35.81 & 35.18 & 37.31 \\
Helmholtz staircase & 0.00 & 0.26 & 0.07 & 0.40 \\
ClimSim & 32.24$\pm$0.06 & 27.20$\pm$0.05 & 26.31$\pm$0.04 & 25.10$\pm$0.18 \\
RT coefficients, state & 71.54$\pm$0.59 & 62.93$\pm$0.51 & 3.97$\pm$0.42 & 3.75$\pm$0.25 \\
RT coefficients, +6S & 33.83$\pm$0.26 & 15.76$\pm$1.98 & 4.60$\pm$0.31 & 2.79$\pm$0.14 \\
\bottomrule
\end{tabular}}
\par\medskip\textbf{C. Ensembles, corrections and combinations}\par\medskip
{\fontsize{9}{11}\selectfont
\setlength{\tabcolsep}{3pt}
\begin{tabular}{lrrrr}
\toprule
corpus & ensemble & \thead{network +\\residual KRR} & feature kernel & stack \\
\midrule
Advection & 11.50$\pm$0.05 & 11.50$\pm$0.05 & 11.64$\pm$0.07 & 11.36$\pm$0.04 \\
Burgers $\nu=0.1$ & 19.28$\pm$0.09 & 2.78$\pm$0.29 & 4.92$\pm$0.30 & 2.78$\pm$0.29 \\
Burgers $\nu=0.01$ & 35.31$\pm$0.70 & 23.12$\pm$0.56 & 29.20$\pm$0.59 & 23.01$\pm$0.54 \\
Burgers $\nu=0.001$ & 38.89$\pm$0.11 & 31.82$\pm$0.11 & 35.59$\pm$0.48 & 31.34$\pm$0.14 \\
Radiative layer & 31.58$\pm$0.02 & 30.40$\pm$0.01 & 31.08$\pm$0.01 & 30.12 \\
Active matter & 36.93 & 35.77 & 36.22 & 35.05 \\
Helmholtz staircase & 0.24 & 0.23 & 0.45 & 0.07 \\
ClimSim & 24.43$\pm$0.17 & 24.33$\pm$0.13 & 24.71$\pm$0.25 & 24.25$\pm$0.17 \\
RT coefficients, state & 3.24$\pm$0.11 & 3.13$\pm$0.05 & 3.52$\pm$0.07 & 3.09$\pm$0.05 \\
RT coefficients, +6S & 2.40$\pm$0.11 & 2.34$\pm$0.09 & 2.80$\pm$0.09 & 2.32$\pm$0.08 \\
\bottomrule
\end{tabular}}

\caption{Ten further configurations: dimensions and run counts (A), reference models (B),
and combinations (C). Errors are mean relative spectral norms in percent, with the
0.05 denominator floor of \eqref{eq:coefficient-floor} for the two RT-coefficient rows.
Spreads are sample standard deviations; the radiative-layer repetitions share one split
and vary initialization, and two rows are single runs. The network budget is 200 epochs
at width 512, except ClimSim (60 at 512) and RT coefficients (120 at 384). The linear
ridge is excluded from the stack; its Helmholtz error is $0.004\%$, rounded in the table.
Kernel solves above 20{,}000 training rows use 6{,}000 Nystr\"om landmarks.}
\label{tab:corpora}
\end{table}

The relative errors vary widely with the problem. Where the map is smooth and the sample small,
Burgers at $\nu = 0.1$, the exact kernel is at $2.81\%$ against the network's $20.34\%$, a factor of
seven, of the same kind as the five- to six-fold gap between the per-input kernel and the network on
the EMIT table's smooth components; as the viscosity falls and the solutions roughen the gap closes
to $23.40\%$ against $36.65\%$ and then to $31.66\%$ against $39.98\%$. Where the sample is large and
the map rough, on ClimSim and on the correction-coefficient corpus at nearly half a million rows, the
network leads instead, $25.10\%$ against $27.20\%$ and $3.75\%$ against $62.93\%$, which is the OCO-2
regime. On advection, where the operator is linear, the linear ridge is the lowest row and every
other family lands within $0.4$ points of it.

Learned metrics help most on the correction-coefficient corpus. From the state alone the isotropic
kernel is at $62.93 \pm 0.51$ and a kernel-flow metric \cite{owhadiyoo} takes it to $3.97 \pm 0.42$,
a factor of sixteen, which brings a kernel on the state within six percent of the network on a
corpus where the plain kernel was seventeen times worse. With the low-fidelity coefficients as
inputs the same pair reads $15.76$ and $4.60$, a factor of $3.4$: the easier the problem is made for
the plain kernel, the less the learned metric adds.

The convex stack is the lowest row on eight of the ten configurations and the second lowest on the
other two; the residual-corrected network is among the lowest two on five. Both exceptions are the
linear problems, and on both the row that beats the stack is the ridge, which the stack does not
contain. On advection the ridge is lower by $0.07$ points while the stack, at $11.36\%$, matches its
best member, the isotropic kernel, to two decimals; that kernel and the learned-metric kernel, which
differ by less than $0.01$ points there, carry between $0.85$ and $0.91$ of its weight at every
split. On the Helmholtz
staircase the ridge reaches $0.004\%$, sixteen times below the next row and about a hundred times
below the weakest member of the pool, and the stack reaches $0.066\%$, slightly below its best member,
the learned-metric kernel at $0.067\%$, on which it puts weight $0.94$. In both cases the stack does
what it is asked within its pool and cannot recover a family the pool excludes. A combination is a
good default when no member of its pool dominates it, and it can be read against the whole
comparison only when the pool contains the families being compared.

These rows use one training budget per corpus and one architecture family, and several of the
corpora have published results from methods tuned for them. With three to five runs, and one on two
rows, the spreads support the large orderings above and not fine distinctions between adjacent
families. Every kernel column also inherits the tuning protocol of Section~\ref{sec:retune}, with
the length scale chosen on a $6{,}000$-row subsample and the winner refit on every row. On the EMIT
table, choosing it with every row in the solve lowers the isotropic kernel's radiance error by a
factor $0.882$ on average, so the kernel columns of Table~\ref{tab:corpora} describe the stated
training and selection procedures rather than the best each kernel can reach.
